\section{Discussion: the Main Differences between the Three Models}
\label{sec:discussion}

\subsection{VAE-MLP vs.\ VAE-Conv: the effect of the network}
The two VAEs share the method and differ only in the network (\cref{tab:arch-compare}).

\paragraph{What the convolutional network improves.}
\begin{itemize}
  \item \textbf{Likelihood:} $5.83$ vs.\ $6.31$\,bits/dim, i.e.\ a $0.49$\,bits/dim better bound.
  \item \textbf{Sample quality:} FID $121.9$ vs.\ $136.0$, and IS $3.88$ vs.\ $3.11$.
  \item \textbf{Coverage:} recall $0.073$ vs.\ $0.003$.
  \item \textbf{Reconstructions:} FID $57.9$ vs.\ $99.4$.
\end{itemize}
Convolutions share filters across positions, and attention relates distant positions. Both give the model the
spatial inductive bias that the MLP, which treats an image as 3072 unrelated numbers, lacks.

\paragraph{What it does not change.} Both VAEs remain in the same regime: blurry, low-recall samples, a large gap
between reconstructions and samples, and a collapsed class distribution. \textbf{A better network improves a VAE
but does not remove its characteristic failure modes.} The paper's own network is not the cause of the VAE's
weaknesses on CIFAR-10. It only makes them worse.

\subsection{VAE-Conv vs.\ DDPM: the effect of the generative approach}
These two models share building blocks, widths, data and optimisation settings (\cref{tab:arch-compare,tab:hparams}),
and their sizes are comparable ($30.4$M vs.\ $35.7$M parameters). Their difference is therefore attributable to
the approach. Following \cref{sec:bg-unify}, three structural differences explain all the measurements:
\begin{enumerate}
  \item \textbf{One factorised step vs.\ many small steps.} The VAE must model the whole image distribution with a
        single Gaussian with independent pixels, which forces blurry averages. That gives low precision, low IS,
        and a realism score of $1.34$ in the blind rating. Each DDPM step only has to model a small denoising
        update, for which a Gaussian is accurate. That gives precision equal to real data, and samples a
        vision-language model cannot tell from real photographs.
  \item \textbf{Learned compressive encoder vs.\ fixed non-compressive encoder.} The VAE's encoder must squeeze the
        image into $128$ numbers and match $\N(\mathbf{0},\I)$ while doing so. It never matches the prior perfectly,
        and that prior hole shows as samples far worse than reconstructions (FID $121.9$ vs.\ $57.9$), collapsed
        recall, and samples that fall between classes. The DDPM's forward process matches its prior by
        construction ($L_T\approx8\times10^{-6}$\,bits/dim) and compresses nothing.
  \item \textbf{Cost of generation.} Both advantages of the DDPM are paid for at sampling time: $1000$ network
        evaluations instead of $1$, so about $1{,}700\times$ slower. DDIM cuts this by up to $100\times$ with no
        retraining, but even 10 steps remain $17\times$ slower than the VAE.
\end{enumerate}

\subsection{Summary}

\begin{table}[H]
\centering\small
\caption{The three models at a glance (best per row in bold).}
\label{tab:summary}
\begin{tabularx}{\textwidth}{@{}l>{\raggedright\arraybackslash}X>{\raggedright\arraybackslash}X>{\raggedright\arraybackslash}X@{}}
\toprule
 & \textbf{VAE-MLP} & \textbf{VAE-Conv} & \textbf{DDPM}\\
\midrule
Sample quality (FID / IS) & 136.0 / 3.11 & 121.9 / 3.88 & \textbf{6.58 / 8.59}\\
Realism / diversity (prec. / rec.) & 0.32 / 0.003 & 0.40 / 0.07 & \textbf{0.69 / 0.58}\\
Likelihood (bits/dim) & 6.31 & 5.83 & \textbf{3.79}\\
Sampling speed (img/s) & \textbf{350{,}102} & 7{,}017 & 4.2 (DDIM-10: 416)\\
Training time (GPU-h) & \textbf{0.44} & 1.32 & 3.94\\
Learned encoder / latent code & yes, 200-d & yes, 128-d & no (latent = noise)\\
Typical failure & blurry colour blobs & blurry, between-class blobs & odd object shapes, sharp texture\\
Best use & teaching baseline & fast sampling, representation learning, compression & high-fidelity generation\\
\bottomrule
\end{tabularx}
\end{table}

\paragraph{Which model to use?} It depends on the requirement:
\begin{itemize}
  \item \textbf{Sample quality, diversity, class coverage or likelihood:} the DDPM is decisively better.
  \item \textbf{Real-time generation, or an encoder for compression or editing:} only the VAE provides these.
\end{itemize}
Modern latent diffusion models~\citep{rombach2022ldm} combine the two: diffusion runs inside the compact latent
space of an autoencoder. Our measurements show why this pairing works. The autoencoder contributes the cheap,
compressed representation the DDPM lacks, and diffusion supplies the sample quality and prior matching the VAE
lacks.

\paragraph{Lessons about evaluation.}
\begin{itemize}
  \item Likelihood and sample quality moved in opposite directions three times: VAE overfitting, the EMA
        ablation, and the DDPM's flat loss while its FID kept falling.
  \item IS alone would hide the VAE's diversity failure, and FID alone would hide the split between realism and
        diversity; precision/recall separates the two.
  \item Far from the data, the metrics stop ranking models consistently: FID and KID disagree on the two VAEs.
\end{itemize}
